\centerline{\bf \Huge Движение I}
\centerline{\bf \Huge Центральная и осевая симметрия}

\begin{flushright}
        \scriptsize{}
\end{flushright}

Отрезок, соединяющий две несоседние вершины многоугольника будем называть диагональю многоугольника.

\problem{1}
Сколько осей симметрии у квадрата?

\problem{2}
Многоугольник будем называть правильным, если все его углы равны и все стороны равны. Сколько осей симметрии есть у правильного $n$---угольника?

\problem{3}
Существует ли многоугольник, у которого ровно две параллельные оси симметрии?

\problem{4}
Из вершины квадрата пущен бильярдный шар под углом $\alpha$. При каком значении угла $\alpha$ шар, отразившись от сторон вернется в эту же вершину?

\problem{5}
Выпуклый многоугольник имеет центр симметрии. Докажите, что сумма градусных мер его углов делится на $360^{\circ}$.

\problem{6}
Докажите, что если фигура имеет две перпендикулярные оси симметрии, то она имеет центр симметрии.

\problem{7}
Даны угол $ABC$ и точка $D$ внутри его. Постройте отрезок с концами на сторонах данного угла, середина которого находилась бы в точке $D$.

\problem{8}
Существует ли шестиугольник, который можно разбить одной прямой на четыре равных треугольника?